% !TEX root = ../main.tex
\corrigesChapitre{3}{Titrages pH-métriques et conductimétriques}
\corrigesSouples

\begin{blocLivre}{cor-3-1}
\begin{corrige}{1}{Une solution commerciale}
\demarche{Raisonner sur un litre de solution, puis utiliser la conservation de la quantité de soluté lors de la dilution.}
\begin{questions}
\item $m_\text{solution}=d\times\rho_\text{eau}\times V=1{,}11\times\qty{1,00}{kg.L^{-1}}\times\qty{1,00}{L}=\res{\qty{1110}{g}}$,
et $m(\ce{HCl})=0{,}23\times1\,110=\res{\qty{255}{g}}$.
\item $c=\dfrac{m(\ce{HCl})}{M\,V}=\dfrac{255}{36{,}5\times1{,}00}=\res{\qty{7,0}{mol.L^{-1}}}$.
\item La quantité de soluté se conserve : $c_0V_0=c_1V_1$, d'où $V_0=\dfrac{0{,}100\times500{,}0}{7{,}0}\approx\res{\qty{7,1}{mL}}$,
prélevés à la pipette graduée et versés dans une fiole jaugée de 500,0 mL contenant déjà de l'eau.
\item Porter lunettes, gants et blouse ; travailler sous hotte ; verser l'acide dans l'eau, et non l'inverse.
\end{questions}
\end{corrige}
\end{blocLivre}
\begin{blocLivre}{cor-3-2}
\begin{corrige}{2}{Exploiter l'équivalence}
\demarche{Écrire l'équation, puis la relation à l'équivalence avec les nombres stœchiométriques.}
\begin{questions}
\item \ce{H3O+ + HO- -> 2 H2O} ; à l'équivalence, $c_\mathrm{A}V_\mathrm{A}=c_\mathrm{B}V_\mathrm{E}$, d'où
$c_\mathrm{A}=\dfrac{0{,}100\times12{,}6}{20{,}0}=\res{\qty{0,0630}{mol.L^{-1}}}$.
\item \ce{NH3 + H3O+ -> NH4+ + H2O} ; $c\,V=c_\mathrm{A}V_\mathrm{E}$, d'où
$c=\dfrac{0{,}0500\times14{,}2}{10{,}0}=\res{\qty{0,0710}{mol.L^{-1}}}$.
\end{questions}
\end{corrige}
\end{blocLivre}
\begin{blocLivre}{cor-3-3}
\begin{corrige}{3}{Lire une courbe de titrage pH-métrique}
\demarche{Tracer deux tangentes parallèles de part et d'autre du saut, puis la droite équidistante.}
\begin{questions}
\item La droite équidistante coupe la courbe au point E : $\res{V_\mathrm{E}=\qty{20,0}{mL}}$ et $\res{\pH_\mathrm{E}=7{,}0}$.
\item $c_\mathrm{A}=\dfrac{c_\mathrm{B}V_\mathrm{E}}{V_\mathrm{A}}=\dfrac{\num{1.00e-2}\times20{,}0}{20{,}0}=\res{\qty{1,00e-2}{mol.L^{-1}}}$.
\item L'acide chlorhydrique est un acide fort : $[\ce{H3O+}]=c_\mathrm{A}$ et $\pH=-\log(\num{1.00e-2})=2{,}0$, valeur
lue sur la courbe à $V_\mathrm{B}=0$ : \res{\text{c'est cohérent}}.
\end{questions}
\schema[corrige]{03-ex03}
\end{corrige}
\end{blocLivre}
\begin{blocLivre}{cor-3-4}
\begin{corrige}{4}{Lire une courbe de titrage conductimétrique}
\demarche{Modéliser chaque portion par une droite ; faire le bilan des ions ajoutés et retirés.}
\begin{questions}
\item Les deux droites se coupent pour $\res{V_\mathrm{E}\approx\qty{7,9}{mL}}$, d'où
$c(\ce{NaOH})=\dfrac{0{,}0500\times7{,}9}{20{,}0}\approx\res{\qty{2,0e-2}{mol.L^{-1}}}$.
\item Avant l'équivalence, chaque ion \ce{H3O+} versé consomme un ion \ce{HO-} ($\lambda=19{,}9$) et apporte un ion
\ce{Cl-} ($\lambda=7{,}63$) : un ion très conducteur est remplacé par un ion moins conducteur, $\sigma$
\res{\text{diminue}}. Après l'équivalence, les ions \ce{H3O+} ($\lambda=35{,}0$) et \ce{Cl-} versés s'accumulent :
$\sigma$ \res{\text{augmente}}, et plus vite qu'elle ne diminuait, car $35{,}0+7{,}63>19{,}9-7{,}63$.
\end{questions}
\schema[corrige]{03-ex04}
\end{corrige}
\end{blocLivre}
\begin{blocLivre}{cor-3-5}
\begin{corrige}{5}{Choisir la verrerie d'un titrage}
\demarche{La solution titrée doit être prélevée avec précision ; la solution titrante doit être versée et mesurée progressivement.}
\begin{questions}
\item La solution titrée se prélève à la \res{\text{pipette jaugée de 20,0 mL}}, la plus précise ; la solution titrante
se verse avec la \res{\text{burette graduée}}, qui mesure le volume versé. L'éprouvette est trop imprécise.
\item Le rinçage élimine l'eau restée dans la burette, qui diluerait la solution titrante et fausserait sa concentration.
\item Non : l'eau ne modifie pas la quantité d'espèce titrée, donc \res{V_\mathrm{E}\text{ ne change pas}} ; seul le pH
(ou la conductivité) des mélanges est modifié.
\end{questions}
\end{corrige}
\end{blocLivre}
\begin{blocLivre}{cor-3-6}
\begin{corrige}{6}{La composition du système au cours d'un titrage}
\demarche{Avant l'équivalence, les ions \ce{HO-} versés sont consommés ; après, ils s'accumulent.}
\begin{questions}
\item \ce{H3O+ + HO- -> 2 H2O} ; $n_0=\num{1.0e-2}\times\num{20.0e-3}=\qty{2,0e-4}{mol}$ et $\res{V_\mathrm{E}=\qty{20,0}{mL}}$.
\item En $10^{-4}$ mol :
\begin{center}\small\renewcommand{\arraystretch}{1.2}
\begin{tabular}{@{}lcccc@{}}\toprule
$V_\mathrm{B}$ & \ce{H3O+} & \ce{Cl-} & \ce{Na+} & \ce{HO-}\\\midrule
10,0 mL & 1,0 & 2,0 & 1,0 & $\approx0$\\
20,0 mL & $\approx0$ & 2,0 & 2,0 & $\approx0$\\
30,0 mL & $\approx0$ & 2,0 & 3,0 & 1,0\\\bottomrule
\end{tabular}\end{center}
\item Pour 10,0 mL : $[\ce{H3O+}]=\frac{\num{1.0e-4}}{\num{30.0e-3}}=\qty{3,3e-3}{mol.L^{-1}}$, $\res{\pH=2{,}48}$.
Pour 30,0 mL : $[\ce{HO-}]=\frac{\num{1.0e-4}}{\num{50.0e-3}}=\qty{2,0e-3}{mol.L^{-1}}$, donc
$[\ce{H3O+}]=\frac{10^{-14}}{\num{2.0e-3}}=\qty{5,0e-12}{mol.L^{-1}}$ et $\res{\pH=11{,}3}$.
\end{questions}
\end{corrige}
\end{blocLivre}
\begin{blocLivre}{cor-3-7}
\begin{corrige}{7}{Interpréter les pentes d'une courbe conductimétrique}
\demarche{Pour une mole d'ions \ce{HO-} versée, sommer les conductivités molaires des ions ajoutés et retrancher celles des ions retirés.}
\begin{questions}
\item \ce{H3O+ + HO- -> 2 H2O} et \ce{CH3COOH + HO- -> CH3COO- + H2O}.
\item Avant l'équivalence : on ajoute \ce{Na+} et on retire \ce{H3O+} : $5{,}01-35{,}0\approx-30$, pente
\res{\text{négative}}. Après : on ajoute \ce{Na+} et \ce{HO-} : $5{,}01+19{,}9=24{,}9$, pente \res{\text{positive}}.
\item Avant l'équivalence : l'acide éthanoïque, très peu dissocié, ne contribue presque pas à $\sigma$ ; on ajoute \ce{Na+}
et il se forme \ce{CH3COO-} : $5{,}01+4{,}09=9{,}10$. La conductivité \res{\text{augmente dès le début}}, lentement.
\item Après l'équivalence, la contribution vaut 24,9 : le rapport des pentes vaut $\frac{24{,}9}{9{,}10}\approx\res{2{,}7}$.
La courbe est formée de deux droites croissantes, la seconde environ 2,7 fois plus raide que la première.
\end{questions}
\schema[corrige]{rev-03-ex07-cor}
\legende{Allure pour l’acide éthanoïque, dilution négligée ; ordonnée en unités arbitraires.}
\end{corrige}
\end{blocLivre}
\begin{blocLivre}{cor-3-8}
\begin{corrige}{8}{Titrer un vinaigre}
\demarche{Déterminer $V_\mathrm{E}$, revenir à la solution non diluée, puis convertir en pourcentage massique.}
\begin{questions}
\item \ce{CH3COOH(aq) + HO^-(aq) -> CH3COO^-(aq) + H2O(l)}.
\item $\res{V_\mathrm{E}=\qty{10,1}{mL}}$ et $\res{\pH_\mathrm{E}\approx8{,}7}$. À l'équivalence, la solution contient des
ions éthanoate, base faible : elle est basique.
\item Solution diluée : $c=\frac{0{,}100\times10{,}1}{10{,}0}=\qty{0,101}{mol.L^{-1}}$ ; vinaigre :
$\res{c=\qty{1,01}{mol.L^{-1}}}$. Un litre (1\,010 g) contient $1{,}01\times60{,}05=\qty{60,7}{g}$ d'acide, soit
$\frac{60{,}7}{1\,010}\times100\approx\res{6{,}0~\%}$ : l'indication est vérifiée.
\item La phénolphtaléine, dont la zone de virage (8,2 -- 10,0) contient le pH à l'équivalence.
\end{questions}
\schema[corrige]{03-ex08}
\end{corrige}
\end{blocLivre}
\begin{blocLivre}{cor-3-9}
\begin{corrige}{9}{Les ions chlorure d'un sérum physiologique}
\demarche{Faire le bilan des ions avant et après l'équivalence, puis tenir compte de la dilution.}
\begin{questions}
\item \ce{Ag^+(aq) + Cl^-(aq) -> AgCl(s)}. Ce n'est pas une réaction acide-base : le pH ne varie pas à l'équivalence.
\item Avant l'équivalence, chaque ion \ce{Ag+} versé précipite un ion \ce{Cl-} ($\lambda=7{,}63$) et apporte un ion
\ce{NO3-} ($\lambda=7{,}14$) : $\sigma$ diminue \res{\text{très légèrement}}. Après l'équivalence, les ions \ce{Ag+}
et \ce{NO3-} s'accumulent : $\sigma$ \res{\text{augmente}}.
\item $\res{V_\mathrm{E}\approx\qty{6,1}{mL}}$ ; solution diluée : $[\ce{Cl-}]=\frac{0{,}0500\times6{,}1}{20{,}0}=\qty{1,5e-2}{mol.L^{-1}}$ ;
sérum : $\res{[\ce{Cl-}]\approx\qty{0,15}{mol.L^{-1}}}$.
\item $c_\mathrm{m}=\num{0.1525}\times58{,}44\approx\qty{8,9}{g.L^{-1}}$ : à la précision de la lecture graphique près,
\res{\text{l'indication de 9,0 g}\cdot\text{L}^{-1}\text{ est vérifiée}}.
\end{questions}
\schema[corrige]{03-ex09}
\end{corrige}
\end{blocLivre}
\begin{blocLivre}{cor-3-10}
\begin{corrige}{10}{Programmer l'évolution des quantités de matière}
\demarche{Une réaction totale consomme entièrement le réactif limitant : on utilise les fonctions \code{max} et \code{min}.}
\begin{questions}
\item \ce{CH3COOH + HO- -> CH3COO- + H2O} ; $n_0=0{,}100\times10{,}0=\res{\qty{1,00}{mmol}}$ (\code{n0 = 1.00}).
\item \code{nA.append(max(n0 - nb, 0))} et \code{nHO.append(max(nb - n0, 0))}.
\item Les ions éthanoate se forment tant qu'il reste de l'acide éthanoïque : leur quantité cesse d'augmenter quand
l'acide est épuisé, à l'équivalence, pour $\res{V_\mathrm{E}=\qty{10,0}{mL}}$.
\item Voir le graphique : la quantité d'ions \ce{Na+}, spectateurs, croît linéairement pendant tout le titrage.
\end{questions}
\schema[corrige]{03-ex10-cor}
\end{corrige}
\end{blocLivre}
\begin{blocLivre}{cor-3-11}
\begin{corrige}{11}{L'acidité d'un lait}
\demarche{Calculer la masse d'acide dans la prise d'essai, puis la rapporter à un litre de lait.}
\begin{questions}
\item $n=c\,V_\mathrm{E}=\frac19\times\num{1.6e-3}=\qty{1,78e-4}{mol}$ et $m=n\,M=\num{1.78e-4}\times90{,}08=\res{\qty{0,0160}{g}}$.
\item Un litre de lait contient $\frac{0{,}0160}{0{,}0100}=\qty{1,60}{g}$ d'acide lactique, soit \res{16~°D} :
inférieur à 18 °D, \res{\text{le lait est frais}}.
\item Pour un volume équivalent $V_\mathrm{E}$ (en mL), la masse d'acide par litre vaut
$\frac{\frac19\times V_\mathrm{E}\times10^{-3}\times90{,}08}{0{,}0100}\approx V_\mathrm{E}\times\qty{1,00}{g.L^{-1}}$,
soit $10\,V_\mathrm{E}$ degrés Dornic : un dixième de millilitre correspond à un degré.
\end{questions}
\end{corrige}
\end{blocLivre}
\begin{blocLivre}{cor-3-12}
\begin{corrige}{12}{Contrôle d'un déboucheur}
\demarche{Exploiter le titrage de la solution diluée, revenir au déboucheur, puis convertir en titre massique.}
\begin{questions}
\item Prélever 10,0 mL de déboucheur à la pipette jaugée (munie d'une propipette), les verser dans une fiole jaugée de
250,0 mL contenant de l'eau, compléter au trait de jauge et homogénéiser. Gants, lunettes et blouse sont indispensables.
\item \ce{HO- + H3O+ -> 2 H2O} ; la méthode des tangentes donne $\res{V_\mathrm{E}=\qty{21,6}{mL}}$.
\item Solution diluée : $c=\frac{0{,}100\times21{,}6}{20{,}0}=\qty{0,108}{mol.L^{-1}}$ ; déboucheur :
$c=25\times0{,}108=\qty{2,70}{mol.L^{-1}}$. Un litre (1\,100 g) contient $2{,}70\times40{,}0=\qty{108}{g}$ de soude :
$t=\frac{108}{1\,100}\approx\res{9{,}8~\%}$. L'écart à 10 \% est inférieur à 2 \% : \res{\text{l'étiquette est vérifiée}}.
\item Pour 10 \% : $c=\frac{0{,}10\times1\,100}{40{,}0}=\qty{2,75}{mol.L^{-1}}$, $c_\text{diluée}=\qty{0,110}{mol.L^{-1}}$ et
$V_\mathrm{E}=\frac{0{,}110\times20{,}0}{0{,}100}=\res{\qty{22,0}{mL}}$.
\end{questions}
\schema[corrige]{03-ex12}
\end{corrige}
\end{blocLivre}
\begin{blocLivre}{cor-3-13}
\begin{corrige}{13}{Titrage d'un comprimé d'aspirine}
\demarche{Repérer $V_\mathrm{E}$ au maximum de la dérivée, puis remonter au comprimé entier.}
\begin{questions}
\item \ce{AH(aq) + HO^-(aq) -> A^-(aq) + H2O(l)}.
\item Le pic de la dérivée, fin et unique, donne $\res{V_\mathrm{E}=\qty{13,8}{mL}}$ ; cette méthode est plus précise que
les tangentes car le saut est bref et la lecture d'un maximum est sans ambiguïté.
\item Dans 50,0 mL : $n=0{,}0500\times\num{13.8e-3}=\qty{6,90e-4}{mol}$ ; dans le comprimé (200,0 mL) :
$4\times\num{6.90e-4}=\qty{2,76e-3}{mol}$, soit $m=\num{2.76e-3}\times180{,}16\approx\res{\qty{497}{mg}}$. L'écart à 500 mg
vaut 0,6 \% : \res{\text{l'indication est vérifiée}}.
\item À chaud, les ions hydroxyde réagissent aussi avec la fonction ester : il en faudrait davantage,
$V_\mathrm{E}$ serait plus grand et la masse d'aspirine \res{\text{surestimée}}.
\end{questions}
\schema[corrige]{03-ex13}
\end{corrige}
\end{blocLivre}
\begin{blocLivre}{cor-3-14}
\begin{corrige}{14}{Un vinaigre frelaté ?}
\demarche{Associer chaque portion de droite à la réaction qui se produit, puis exploiter les deux volumes équivalents.}
\begin{questions}
\item D'abord \ce{H3O+ + HO- -> 2 H2O}, puis \ce{CH3COOH + HO- -> CH3COO- + H2O}.
\item Portion 1 : les ions \ce{H3O+} (35,0) sont remplacés par des ions \ce{Na+} (5,01) : forte diminution.
Portion 2 : apparition d'ions \ce{Na+} et \ce{CH3COO-} ($5{,}01+4{,}09$) : faible augmentation.
Portion 3 : accumulation d'ions \ce{Na+} et \ce{HO-} ($5{,}01+19{,}9$) : forte augmentation.
\item $V_\mathrm{E1}\approx\qty{4,1}{mL}$ : $c(\ce{HCl})=\frac{0{,}100\times4{,}1}{20{,}0}\approx\qty{0,020}{mol.L^{-1}}$ dans S ;
$V_\mathrm{E2}-V_\mathrm{E1}\approx\qty{9,8}{mL}$ : $c(\ce{CH3COOH})=\frac{0{,}100\times9{,}8}{20{,}0}=\qty{0,049}{mol.L^{-1}}$ dans S.
Dans le vinaigre, vingt fois plus concentré : $\res{c(\ce{CH3COOH})\approx\qty{0,98}{mol.L^{-1}}}$ et
$\res{c(\ce{HCl})\approx\qty{0,41}{mol.L^{-1}}}$. La présence d'acide chlorhydrique montre que
\res{\text{le vinaigre est frelaté}}.
\end{questions}
\schema[corrige]{03-ex14}
\end{corrige}
\end{blocLivre}
\begin{blocLivre}{cor-3-15}
\begin{corrige}{15}{Quand la dilution n'est plus négligeable}
\demarche{La loi de Kohlrausch s'écrit avec les quantités de matière divisées par le volume total.}
\begin{questions}
\item $\sigma=\sum\lambda_i\dfrac{n_i}{V_0+V_\mathrm{B}}$, donc $\sigma\,(V_0+V_\mathrm{B})=\sum\lambda_in_i$. De part et
d'autre de l'équivalence, chaque $n_i$ est une fonction affine de $V_\mathrm{B}$ : \res{\sigma\,(V_0+V_\mathrm{B})\text{ est affine}}.
\item $\sigma'=4{,}21\times\frac{20{,}0+10{,}0}{20{,}0}=\res{\qty{6,32}{mS.cm^{-1}}}$.
\item $c=\frac{0{,}100\times10{,}0}{20{,}0}=\res{\qty{0,0500}{mol.L^{-1}}}$.
\item L'erreur relative vaut $\frac{10{,}0-7{,}9}{10{,}0}\approx\res{21~\%}$ : le volume équivalent est fortement
sous-estimé. On évite le problème en ajoutant un grand volume d'eau (par exemple 200 mL) avant le titrage :
la dilution due à la solution titrante devient négligeable et les points s'alignent.
\end{questions}
\schema[corrige]{03-ex15}
\end{corrige}
\end{blocLivre}
